% Preamble for Reference Grammar PDF build (pandoc + tectonic).
% Loaded via --include-in-header.

% --- Fonts ----------------------------------------------------------------
% Brill v4 (non-commercial license per assets/fonts/Brill_typeface_end_user_license_agreement.pdf).
% Brill is a scholarly serif with comprehensive IPA, Greek, and combining-diacritical
% coverage, designed for humanities use. It replaces the v1 setup of EB Garamond +
% Charis SIL fallback chain.
\usepackage{fontspec}
\setmainfont{Brill}[
  Path           = assets/fonts/,
  Extension      = .ttf,
  UprightFont    = *-Roman,
  ItalicFont     = *-Italic,
  BoldFont       = *-Bold,
  BoldItalicFont = *-BoldItalic,
  Ligatures      = TeX
]

% DejaVu Sans Mono as the monospace face. Latin Modern Mono (the default)
% lacks Greek/IPA, which we need for the small ASCII metrical-tree code
% block and a handful of inline backticked glyphs.
\setmonofont{DejaVuSansMono}[
  Path           = assets/fonts/,
  Extension      = .ttf,
  UprightFont    = *,
  ItalicFont     = *-Oblique,
  BoldFont       = *-Bold,
  BoldItalicFont = *-BoldOblique,
  Scale          = 0.85
]

% Note on IPA / combining-diacritical / Greek coverage: Brill provides these
% natively, so the Charis SIL fallback chain (XeTeX intercharclass routing
% IPA blocks to Charis) that the EB Garamond setup required has been removed.
% The build pre-pass in build.ps1 still rewrites U+2208 (∈) to inline math,
% kept defensively because longtable cell context interferes with font
% selection for math-symbol blocks even when the main font has the glyph.

% Route all monospace (backtick spans) through Brill — no DejaVu in output.
\let\ttfamily\rmfamily

% --- Typography polish ----------------------------------------------------
\usepackage{microtype}
\usepackage{setspace}
\setstretch{1.14}

\setlength{\parskip}{0.45em plus 0.1em minus 0.05em}
\setlength{\parindent}{0pt}

% --- Heading styles -------------------------------------------------------
\usepackage{titlesec}

\titleformat{\chapter}[display]
  {\normalfont\sffamily\filright}
  {\large\textsc{Chapter \thechapter}}
  {0.5em}
  {\Huge\bfseries\rmfamily}
\titlespacing*{\chapter}{0pt}{-20pt}{24pt}

\titleformat{\section}
  {\normalfont\Large\bfseries}{\thesection}{0.7em}{}
\titleformat{\subsection}
  {\normalfont\large\bfseries}{\thesubsection}{0.6em}{}
\titleformat{\subsubsection}
  {\normalfont\normalsize\bfseries\itshape}{\thesubsubsection}{0.6em}{}

% --- Headers & footers ----------------------------------------------------
\usepackage{fancyhdr}
\setlength{\headheight}{14pt}
\pagestyle{fancy}
\fancyhf{}
\fancyhead[L]{\itshape\nouppercase{\leftmark}}
\fancyhead[R]{\thepage}
\renewcommand{\headrulewidth}{0pt}

\fancypagestyle{plain}{%
  \fancyhf{}%
  \fancyfoot[C]{\thepage}%
  \renewcommand{\headrulewidth}{0pt}%
}

% --- Tables ---------------------------------------------------------------
\usepackage{booktabs}

% Suppress pandoc's default \maketitle so the custom titlepage owns the cover.
% (Hyperref colors are handled via -V flags in build.ps1, since pandoc loads
% hyperref *after* this header-include block runs.)
\let\oldmaketitle\maketitle
\renewcommand{\maketitle}{}
